\section{Introduction and current theorem structure}\label{sec:current-overview}\label{sec:introduction}
A finite-dimensional physical state can generate responses whose exact classical realization requires arbitrarily large memory. The question is not answered by finite linear rank: positivity, a legal stochastic decoder and the prescribed initial state all matter. We determine this question for the balanced single-frequency response of a damped qubit, and connect the answer to the actual least-cost optimum of the fixed squared task in Theorem~\ref{thm:r67-embedding}.

Theorem~\ref{thm:r69-spectral} identifies the exact live-label count with a classical stochastic spectral region, with no additional initial label. The proof uses the three nonzero response modes in one direction and a maximal-modulus eigenvector coordinate in the other. This step retains arbitrary phase and every nonzero contrast. It does not give a decoder with access to the protected system, nor one decoder for an unknown family of initial phases.

The central quantitative consequence is arithmetic. Theorem~\ref{thm:r69-worst-phase} gives the exact worst-phase size and its square-root divergence as damping disappears. Theorem~\ref{thm:r69-cubic} proves a matching cubic-root divergence for each fixed badly approximable phase, with explicit constants and a global elementary example. A rational approximation selects a small invariant polygon for the upper; applying a stochastic sector inequality to a suitable power of every competing realization gives the lower. Theorem~\ref{thm:r69-finite-mixture} extends the matching exponent to finite spectral mixtures on higher-dimensional protected carriers, counting the initial mixture label. Both bounds refer to exact all-time responses at the same least expected-cost budget. The law is therefore not inferred from finite-time simulation, a fitted scaling curve or a numerical eigenvalue plot.

The underlying stochastic eigenvalue regions, sector bound and first-sector lazy-cycle rigidity are classical \cite{DmitrievDynkin46,JohnsonPaparella17,KirklandSmigoc22}. We prove the needed sector statement directly rather than use an unproved external boundary theorem. Noise-dependent classical-memory lower bounds for quantum processes also precede this work \cite{Fanizza24}. The distinction claimed here is the exact prescribed-start response equivalence, the matched fixed-phase arithmetic law and their implementation in the common squared-loss experiment. No exhaustive priority or journal-significance certificate is inferred from these comparisons.

\subsection*{Other current structural results and their separate hypotheses}
A finite physical system, a finite classical model of its observations, and a finite observer implementing an optimal decision are different mathematical objects. The central question of General Theta Foundations I is how a prepared causal experiment turns physically acquired information into attainable prediction under explicit resources. The actual-history transfer theorem compares task risk with acquired geometry under its paid converse and common-construction certificates. The performed-law theory establishes unrestricted attainment, and the boundary results determine when exact observer memory follows. We now supply a constructive reduction of a broad class of controlled physical instruments without conflating any of these implications.

Theorem~\ref{thm:r68-instrument} replaces a supplied invariant-cone certificate by a check on each actual command/report branch. If all nonzero normalized outputs have a common positive spectral floor, one pure-state net works for every branch, even when commands do not commute or share an invariant state. The resulting hidden classical instrument preserves the complete observations of every adaptive observer and every terminal measurement. Its expected-cost identity is proved at each actual begun history and summed monotonically; it is not inferred from a bounded-loss discrepancy. A block extension counts all intermediate reports and its clock rather than making block calls free.

Theorem~\ref{thm:r68-noise-exponent} gives matching powers of the mixture margin for uniform state-positive realization at each fixed physical dimension. The exponent comes from the geometry of pure-state projective caps. Mixed hidden density vertices cannot evade that bound. Finite physical command families already force the same exponent, although their command count grows. This is deliberately not a universal lower bound for all hidden response representations: a classical generator need not correspond to a quantum density matrix when only one task is prescribed. The separate fixed-angle spectral bound, Corollary~\ref{cor:r68-qubit-lower}, applies to all stochastic scalar-response realizations and is explicitly non-sharp.

Positive realization, invariant polyhedral cones and finite-rank obstructions are established ideas \cite{Heller65,CarlylePaz71,MonrasWinter14}. Fanizza, Lumbreras and Winter also obtain noise-dependent classical-memory obstructions for quantum processes \cite{Fanizza24}. The contribution asserted here is the explicit branchwise construction, its dimension-dependent uniform count and sparse rows, the exact adaptive stopped-resource identity, and the distinctions between its resource classes. We do not claim to originate the cone mechanism or the classical--quantum separation.

The previously proved single-command boundary criterion and squared-audit hierarchy remain in the main article below, along with the complete actual-history comparison proof. The robust optimal face includes global modelwise inequalities even when every Bayes action is tied. Neither a finite physical hidden model nor a finite linear Hankel rank alone settles the exact implementation of that face. The full performed-law, primal--dual, timing, finite-width, reset, spectral, geometry and predictive-quotient developments remain current in the paired supplement; all predecessor sources remain immutable ancestry.
